\section{Simulation Model}
\label{sec:simulation}

The simulation runs in Gazebo Fortress with ROS~2 Humble, inside a Docker container that is
not privileged and receives only the graphics card, the display socket and the joystick
devices of the host. One script builds the image, starts the container and launches every
configuration described below.

\subsection{From CAD to a Robot Description}
\label{sec:urdf}
The CAD assembly was exported from SolidWorks as a URDF description with one mesh per link.
An export of this kind is a starting point and not a usable model. The description used
here is maintained by hand on top of the export, and it differs from it in the following
ways.
\begin{itemize}
  \item A frame \code{base\_footprint} on the floor was added below \code{base\_link}, which
        stays at the midpoint of the axle, as REP~105 prescribes~\cite{rep105}.
  \item The wheel axes were turned so that a positive rotation drives the robot forward.
  \item The right wheel was made coaxial with the left one. The export had placed it
        0.5\unit{mm} higher.
  \item The limits, effort, velocity and damping of the gripper joints, all zero in the
        export, were filled in from the data of the servo.
  \item The body of the depth camera, which was not part of the CAD assembly, was added as a
        box of $90\times25\times25\unit{mm}$ and 72\unit{g}, together with the depth and
        optical frames that the camera driver publishes.
  \item The infrared array was given a mass. The export had assigned none, and the simulator
        rejects a link without mass.
\end{itemize}

\paragraph{Collision geometry.} The visual mesh of the chassis has 277\,864 triangles. For
contact it is replaced by primitives: a box for the plate, a box for the volume above it, a
cylinder for each wheel and a sphere for the caster. The jaws and the camera holder keep
their meshes, of about 2\,000 and 700 triangles, because the shape of the jaws is what the
gripper depends on.

\paragraph{The caster.} In the CAD model the bottom of the caster dome is 0.74\unit{mm}
above the plane on which the wheels stand, so the robot would rock on it. The collision
sphere of the model is sized to touch the floor exactly. The description records this as a
correction to be made on the physical part.

\subsection{Drive and Contact}
\label{sec:sim-drive}
The wheels are driven by the differential-drive system of Gazebo, which converts a velocity
command into wheel speeds by \cref{eq:diffdrive} and publishes wheel odometry at
30\unit{Hz}. The physics step is 1\unit{ms}. The friction coefficients are 1.0 for the
wheels, 0 for the caster, 0.8 for the jaws and 0.35 for the coins.

The drive is limited to 1.0\unit{m/s^2} and 8\unit{rad/s^2}, in acceleration and in
braking. The braking limit was added after the first runs: without it the wheels stopped
at once, and the robot skidded and turned. \Cref{sec:mass}
gives the reason, a wheel load of 18\,\% of the weight, and the limit of
\cref{eq:traction} that the chosen value respects.

\subsection{Gripper}
The two jaw joints are position-controlled at the rated speed of the servo. Gazebo Fortress
ignores the mimic relation by which a URDF couples one joint to another, so the gear pair is
reproduced in software: a node receives the commanded angle of the driven jaw, clamps it to
the joint limits, rejects values that are not finite, and sends the angle to one jaw and its
negative to the other.

\subsection{Sensor Models}
\label{sec:sensor-models}
\paragraph{Depth camera.} The camera is simulated as an RGB-D sensor at the position of the
left infrared imager of the D435, with the depth field of view of the real camera,
$87^\circ\times58^\circ$, an image of $848\times480$ pixels, a depth range from 0.2 to
10\unit{m} and Gaussian image noise with a standard deviation of 0.007. The topic names
match those of the driver of the real camera, so that the same nodes can run on the
hardware. The point cloud is computed in ROS from the depth image, because the cloud that
the simulator produces is expressed in a different axis convention under the same frame
name.

\paragraph{Infrared array.} Gazebo has no reflectance sensor. Each channel is therefore
modelled by a camera of $8\times8$ pixels that looks straight down from the sensing face
and covers a spot about 8\unit{mm} across on the floor, comparable to the spot of a
reflectance sensor at that height. A node averages the brightness of each small image,
scales it to a voltage between 0 and 5\unit{V}, adds Gaussian noise, and publishes the eight
voltages at 20\unit{Hz}. A channel that lies over the edge of a line therefore returns an
intermediate voltage, as the real sensor does, and it is this that allows the position of a
line to be interpolated between channels (\cref{sec:linefollower}). The arena is lit evenly
and without shadows, so that the readings do not depend on the heading of the robot. In the
mission runs the noise has a standard deviation of 0.05\unit{V}.

\paragraph{Inertial sensor.} The inertial sensor is sampled at 100\unit{Hz}, with Gaussian
noise of 0.002\unit{rad/s} on the angular rate and 0.02\unit{m/s^2} on the acceleration.

\subsection{What the Model Leaves Out}
The motors are ideal velocity sources within the acceleration limits: there is no model of
the gear train, of backlash or of the motor driver. The servo is an ideal position source.
The encoders have no quantisation. The infrared model has no dependence on ambient light or
on the height of the array above the floor. The floor is flat. These omissions bound what
the results of \cref{sec:results} can say about the hardware, and \cref{sec:discussion}
returns to them.
